%% =====================================================================
\section{The scan}\label{sec:scan}
%% =====================================================================

With the ledger validated and the classification in hand, the remaining question is
empirical: does anything in the family beat what is known, and if not, exactly what stops
it? This section reports three computations --- an optimality computation, a rank-reduction
search, and a coset-cancellation sweep --- each with its domain, its tolerance and, where
applicable, its positive control.

\subsection{Four doors that had to be re-opened}\label{ss:fourdoors}

A first optimiser over \eqref{eq:exported} necessarily leaves gaps, and it is worth naming
them, because two of the four changed the answer.
\begin{enumerate}[label=\textup{(\arabic*)},leftmargin=2.6em,itemsep=3pt]
\item $m=-1$, the primitive convention, is not a valid slice: Finding~\ref{find:mminus1}.
\item $E$ at $m\ge2$ had never been measured, every anchor having $m\le1$:
      Finding~\ref{find:R1R2}. This changed the $\zeta_2(7)$ answer by $4$ units.
\item The \emph{rank} axis. Reporting only rank-$1$ cells hides the fact that at $p=2$ the
      size constraint is never binding: \S\ref{ss:parity}.
\item The $O(1)$ \emph{inset} cone --- bounded translations of the bricks. These move no
      asymptotic rate at all, so they are invisible to the ledger, but they move every
      rational coefficient, and so they are exactly where a rank reduction could hide:
      \S\ref{ss:rank}.
\end{enumerate}

Before anything else, the scan code reproduces the anchors of \S\ref{sec:validation}
independently: LSZ's printed $\rho_{0,3}=768$, $\rho_{1,3}=73728$, $\rho_{0,0}=0$,
$\rho_{1,0}=-1024$, the normalisations $\rho_n=1,96,14944$, and all three published measures
to every digit \VER{exact}.

\subsection{The optimisation domain}\label{ss:optdomain}

We state the domain first, so that the theorem that follows is unambiguous.

\begin{definition}[the cone]\label{def:cone}
The \emph{cone} $\mathcal C$ consists of the rational functions
\begin{equation}\label{eq:conefam}
 R(t)=(2t+h_0)^{\epsilon}\,
 \frac{\prod_{j\le M}(t+\theta'_j+e_j)_{h_0-2e_j}}{\prod_{j\le M}(t+f_j)_{h_0+1-2f_j}}
\end{equation}
obtained from Definition~\ref{def:family} by allowing, in addition to the parameter tuple
$(p,r,\text{shifts},A,m,\epsilon)$ with
\[
 p\le13,\quad r\le3,\quad A\le14,\quad 0\le m\le w+2,\quad \epsilon\in\{0,1\},
 \quad\text{all coset spreads},\quad w\le9,
\]
the following brick displacements:
\begin{enumerate}[label=\textup{(\alph*)},leftmargin=2.6em,itemsep=2pt]
\item \emph{$O(1)$ insets}: integer $e_j,f_j\ge0$ bounded independently of $n$ ($58$, $127$
      and $288$ admissible points at $M=3,4,5$ respectively);
\item \emph{scaled insets}: $e_j=\eta_jn$, $f_j=\phi_jn$ with $\eta_j,\phi_j\ge0$ --- this is
      the Rhin--Viola length cone, and it is where $\dgr>0$ lives;
\item \emph{the group orbit}: the whole $1920$-element $\Gp$-orbit of the direction, via
      Corollary~\ref{cor:transfer};
\item \emph{rank reduction}: all $\QQ$-linear combinations
      $L=\sum_{j,l,d}c_{j,l,d}\,n^d\,(n\to n+l)\,(\text{inset-form }j)$ with polynomial
      coefficients, within the shift/degree bounds of \S\ref{ss:rank}.
\end{enumerate}
Negative pole insets are excluded throughout, since they put poles on the wrong side of
$\Zp$.
\end{definition}

\begin{remark}[no calibration step]\label{rem:nocalib}
In the computations of this section $E$ and the archimedean growth are \emph{measured}, not
taken from \eqref{eq:E} and \eqref{eq:vpC}. Consequently any $\Phi_n$-type or Rhin--Viola
saving that is actually present at a point of $\mathcal C$ is already counted in the number
reported for that point. There is no fitting step anywhere in \S\ref{sec:scan}.
\end{remark}

\subsection{Optimality of the three known measures}\label{ss:optimal}

\begin{theorem}\label{thm:optimal}
\COMP{the cone $\mathcal C$ of Definition~\ref{def:cone}} Over $\mathcal C$, the three
published $p$-adic irrationality measures are attained and cannot be improved:
\[
\begin{aligned}
 \mu(\zeta_2(3))&\le7.17739889912418, &\quad
 \mu(\zeta_3(3))&\le22.28144795149432,\\
 \mu(\zeta_2(5))&\le20.34265173891448, &&
\end{aligned}
\]
and no configuration in $\mathcal C$ yields a smaller bound for any of the three, nor a
finite bound for any other $\zeta_p(w)$ with $p\le13$, $w\le9$.
\end{theorem}

We emphasise, as in \S\ref{ss:scope}: this is optimality \emph{over $\mathcal C$}. It is not
a statement that these bounds are unimprovable.

The evaluation is of $\mu=\alpha/(\alpha-\beta)=(G+\mathrm{gain})/(\mathrm{gain}-C_1-E)$
with $E$ and the growth measured. The four sub-computations:

\begin{enumerate}[label=\textup{(\roman*)},leftmargin=2.6em,itemsep=4pt]
\item \textbf{$O(1)$ insets: asymptotically no effect, and never an improvement.} At all
      $58/127/288$ admissible points ($M=3,4,5$) every \emph{asymptotic} quantity ---
      $\sum\lambda=\sum\nu=M$, $G$, the gain, $C_1=0$ --- is identical to the symmetric
      point, and the measured $E$ is identical too, so $\mu$ is asymptotically unchanged,
      exactly. At finite $n$ the insets perturb $\sum\lambda,\sum\nu$ by $O(1/n)$, and the
      symmetric point $e=f=0$ is the best point of the cone at every $n$ tested, at both
      rank-$1$ points; at the rank-$2$ control $M=5$ it is the best point at $n\ge36$ but
      not at $n=24$, for the artefactual reason recorded next.
%% REFEREE [R15]: "the best point of the cone at every n tested" is contradicted at
%% n = 24 by the M=5 block, where s_record.py reports BEST at e=f=(0,0,2,2,2) with E=3 --
%% the very artefact recorded two sentences later.  Proviso moved in front of the claim.
      \VER{$473$ points; exact; $n=24$, spot-checked at $n=36,48,60$}
      \emph{A caveat honestly recorded}: a promising drop $E:5\to3$ at $M=5$, $n=24$ was
      chased and is a small-$n$ artefact; at $n=36,48,60$ the exponent is back to $5$.
      See Appendix~\ref{app:repro} for a second finite-$n$ artefact in the same printout.

\item \textbf{Scaled insets: strictly downhill.} \S\ref{ss:cone}: $G$ falls and $E$ rises
      simultaneously, so $\mu$ is worse everywhere off the symmetric point.

\item \textbf{Rank reduction: blocked.} \S\ref{ss:rank}. This is the only place where a
      \emph{record} was available on paper --- $\mu(\zeta_2(5))\le7.17740$ from the $M=5$,
      $m=0$ family and $\mu(\zeta_2(7))\le12.62404$ from $M=6$, $m=1$ --- and both are closed.

\item \textbf{The group: $\dgr\le0$ on the $p$-adic locus.} Proposition~\ref{prop:symzero}
      gives $\dgr=0$ at the totally symmetric point \PROVED{}. The scan adds the empirical
      half: on the whole $p$-adic locus the \emph{measured} denominator never falls below
      its symmetric value, so $\dgr\le0$ throughout \VER{the cone of \S\ref{ss:cone}}. The
      disjoint-support wall of \S\ref{ss:tension} is thereby confirmed at $p=2$ as well as at
      $p\ge5$.
\end{enumerate}

Re-verified at $50$ decimal digits from $(p,r,A,m,\epsilon)$ alone:
{\footnotesize
\[
\begin{array}{llll}
 \zeta_2(3): & G=4.158883083359672, & \marg=1.158883083359672, & \mu=7.1773988991241796616,\\
 \zeta_3(3): & G=3.295836866004329, & \marg=0.2958368660043291, & \mu=22.28144795149432156,\\
 \zeta_2(5): & G=5.545177444479562, & \marg=0.5451774444795625, & \mu=20.34265173891448181 .
\end{array}
\]
}

\subsection{The parity law, and why size never binds at \texorpdfstring{$p=2$}{p=2}}
\label{ss:parity}

\begin{proposition}[parity law]\label{prop:parity}
\PROVED{}, and \VER{$M=3,\dots,7$, exact} For the very-well-poised family
\eqref{eq:conefam} with any integer insets $e,f\ge0$ one has
$R(-h_0-t)=-(-1)^MR(t)$, hence $r_{i,h_0-k}=-(-1)^{M+i}r_{i,k}$ and therefore
\[
 \rho_i=0\qquad\text{whenever }i+M\text{ is even.}
\]
\end{proposition}

\begin{corollary}\label{cor:collapse}
\DERIVED{} Combining Proposition~\ref{prop:parity} with $\rho_1=0$
(Lemma~\ref{lem:rho1}) and $\zeta_2(\text{even})=0$ (K2), the entire ledger of the $p=2$
very-well-poised family collapses to two numbers:
\begin{equation}\label{eq:MMlaw}
 \marg(M,m)=2M\log2-(M+m),\qquad
 \text{weights }\ \begin{cases}\{m+4,m+6,\dots,m+M\}, & M\ \text{even},\\
                               \{m+3,m+5,\dots,m+M\}, & M\ \text{odd},\end{cases}
\end{equation}
with $\rk$ equal to the number of weights listed.
\end{corollary}

Since $\marg(M,m)=M(2\log2-1)-m\to+\infty$ as $M\to\infty$, we obtain:

\begin{finding}[size is not the problem at $p=2$; rank is]\label{find:sizenotproblem}
\COMP{$\mathcal C$ at $p=2$} There is a positive margin at \emph{every} weight at $p=2$:
\begin{center}
\tabsmall
\begin{tabular}{cclcccll}
\toprule
$M$ & $m$ & weights & $\rk$ & $G$ & $E$ & $\marg$ & $\mu$\\
\midrule
$3$ & $0$ & $[3]$ & $1$ & $4.1589$ & $3$ & $+1.1589$ & $7.17740$ \ (Beukers, $\zeta_2(3)$)\\
$4$ & $1$ & $[5]$ & $1$ & $5.5452$ & $5$ & $+0.5452$ & $20.34265$ \ (LSZ, $\zeta_2(5)$)\\
$4$ & $3$ & $[7]$ & $1$ & $5.5452$ & $11$ & $-5.4548$ & --- \ (the rank-$1$ $\zeta_2(7)$ slice)\\
$6$ & $1$ & $[5,7]$ & $2$ & $8.3178$ & $7$ & $+1.3178$ & $12.62404$ \ (``one of $\zeta_2(5),\zeta_2(7)$'')\\
$7$ & $0$ & $[3,5,7]$ & $3$ & $9.7041$ & $7$ & $+2.7041$ & $7.17740$\\
$8$ & $1$ & $[5,7,9]$ & $3$ & $11.0904$ & $9$ & $+2.0904$ & $10.61098$\\
$9$ & $0$ & $[3,5,7,9]$ & $4$ & $12.4766$ & $9$ & $+3.4766$ & $7.17740$\\
\bottomrule
\end{tabular}
\end{center}
\emph{The entire obstruction at $p=2$ is that the form carries $\lfloor(M-1)/2\rfloor$ zeta
values.} If the $\zeta_2(5)$-coefficient of the $M=6$, $m=1$ family could be removed,
$\zeta_2(7)$ would be irrational with $\mu\le12.62404$.
\end{finding}

Corollary~\ref{cor:collapse} also produces the first of the two $\zeta_2(7)$ walls. Rank $1$
forces $M\in\{2,3\}$ ($m$ even) or $M\in\{3,4\}$ ($m$ odd), hence $M\le4$; then weight $7$
needs $m=3$ ($M=4$) or $m=4$ ($M=3$), and in both cases $m\ge A-1$, so rule (R2) of
Finding~\ref{find:R1R2} bites and the second prime band appears.

\begin{center}\fbox{\begin{minipage}{0.92\textwidth}
\textbf{Wall \#1.} Every rank-$1$ route to weight $7$ at $p=2$ must buy the weight with
derivatives, and every such derivative order switches the second denominator band on. This
is structural, not numerical.
\end{minipage}}\end{center}

\subsection{The rank wall}\label{ss:rank}

Removing an unwanted coefficient from a rank-$R$ form requires an operator
\begin{equation}\label{eq:operator}
 L=\sum_{j,l,d}c_{j,l,d}\;n^d\;(n\to n+l)\;(\text{inset-form }j)
\end{equation}
with \emph{polynomial} coefficients. Anything of size $e^{Gn}$ doubles $\beta$ and gives
$\alpha-2\beta=-2E<0$; this is the classical elimination tax, and it is not a modelling
assumption: combining the $M=6$ form with the LSZ form explicitly gives margin
$11.09-25.86=-14.78$. So the question is precisely whether
$\ker(\rho_{w'})\subseteq\ker(\rho_w)$ inside the difference module generated by the family.

\begin{finding}[five elimination searches, $0/300$]\label{find:rankwall}
\EXCL{exact linear algebra over $\FF_q$, $q=2^{61}-1$; zero tolerance; degree and shift
bounds as tabulated; every hit re-verified over $\QQ$}
\begin{center}
\tabsmall
\begin{tabular}{llcccccccc}
\toprule
test & family & kill $\to$ keep & forms & $L$ & $D$ & unknowns & $\dim\ker$ & keep $\ne0$ & control\\
\midrule
shifts only & $M{=}6$, $m{=}1$ & $\rho_3\to\rho_5$ & $1$ & $4$ & $10$ & $55$ & $0$ & --- & ---\\
shifts, wider & $M{=}6$, $m{=}1$ & $\rho_3\to\rho_5$ & $1$ & $6$ & $12$ & $91$ & $0$ & --- & ---\\
inset cone & $M{=}6$, $m{=}1$ & $\rho_3\to\rho_5$ & $4$ & $3$ & $8$ & $144$ & $66$ & $\mathbf 0$ & $66/66$\\
inset cone, wider & $M{=}6$, $m{=}1$ & $\rho_3\to\rho_5$ & $6$ & $4$ & $6$ & $210$ & $120$ & $\mathbf 0$ & $120/120$\\
the $\zeta_2(5)$ record route & $M{=}5$, $m{=}0$ & $\rho_2\to\rho_4$ & $4$ & $3$ & $8$ & $144$ & $77$ & $\mathbf 0$ & $77/77$\\
the rank-$3$ route & $M{=}7$, $m{=}0$ & $\rho_2,\rho_4\to\rho_6$ & $4$ & $3$ & $8$ & $144$ & $36$ & $\mathbf 0$ & $36/36$\\
the rank-$3$ route & $M{=}8$, $m{=}1$ & $\rho_3,\rho_7\to\rho_5$ & $4$ & $3$ & $6$ & $112$ & $1$ & $\mathbf 0$ & $1/1$\\
\bottomrule
\end{tabular}
\end{center}
Totalling the five searches with a non-trivial kernel: $66+120+77+36+1=300$ kernel vectors,
of which $\mathbf{0}$ keep the wanted coefficient, against positive controls at
$\mathbf{300/300}$.
\end{finding}

\begin{remark}[the positive control]\label{rem:control}
The control evaluates the \emph{same} kernel on a coefficient sequence taken from a
\emph{different} family ($M=8$, resp.\ $M=6$). Every kernel vector gives a nonzero value
there. So the machinery does detect survivors when survivors exist, and the negative result
is about the mathematics and not about the search. The two ``shifts only'' rows have
$\dim\ker=0$ and hence carry no control; they are reported for completeness.
\end{remark}

\begin{center}\fbox{\begin{minipage}{0.92\textwidth}
\textbf{Wall \#2.} Every polynomial-coefficient operator that annihilates the $\zeta_2(5)$
coefficient annihilates the $\zeta_2(7)$ coefficient as well --- and the same in all four
families tested, including the $\zeta_2(5)$-record route. The coefficients generate the same
difference module; these forms are irreducible.
\end{minipage}}\end{center}

The statement that would upgrade Wall \#2 to a theorem is Open Lemma~\ref{ol:kernel} below.

\subsection{The length cone at \texorpdfstring{$p=2$}{p=2}: the tension, measured}\label{ss:cone}

The remaining freedom in $\mathcal C$ is the length cone --- $e_j=\eta_jn$, $f_j=\phi_jn$,
the Rhin--Viola direction. With $e,f\ge0$ forced,
\[
 \textstyle\sum\lambda=M-2\sum\eta\le M,\qquad \sum\nu=M-2\sum\phi\le M,
 \qquad G=(\sum\lambda+\sum\nu)\log2\le2M\log2,
\]
with equality \emph{exactly} at the symmetric point --- where $C_1=0$ and $\dgr=0$.

\begin{finding}\label{find:cone}
\VER{$M=6$, $m=1$; $C_1$ de-biased against the symmetric control at the same $h_0$; $E$
measured} Along the length cone:
\begin{center}
\tabsmall
\begin{tabular}{lcccccc}
\toprule
profile $(\eta;\phi)$ & $\sum\lambda$ & $\sum\nu$ & $G$ & $C_1$ & $E$ & $\marg$\\
\midrule
symmetric & $6.000$ & $6.000$ & $8.3178$ & $+0.0000$ & $7$ & $+1.3178$\\
one numerator brick shortened & $5.667$ & $6.000$ & $8.0867$ & $-0.6408$ & $9$ & $-0.9133$\\
two numerator bricks shortened & $5.333$ & $6.000$ & $7.8557$ & $-1.1874$ & $10$ & $-2.1443$\\
one numerator $+$ one pole & $5.667$ & $5.667$ & $7.8557$ & $-0.0260$ & $7$ & $+0.8557$\\
stronger asymmetry & $5.000$ & $5.667$ & $7.3936$ & $-1.3089$ & $9$ & $-1.6064$\\
\bottomrule
\end{tabular}
\end{center}
$G$ falls \emph{and} $E$ rises. The symmetric point is strictly optimal.
\end{finding}

Because $E$ is measured here, this already includes whatever $\Phi_n$ saving is actually
present. \emph{At $p=2$ the $C_1/\dgr$ tension is not a near miss: $\dgr$ is $0$ and the cone
is strictly downhill.}

\subsection{\texorpdfstring{$\zeta_2(7)$}{zeta 2(7)}: the verdict}\label{ss:z27}

\begin{finding}\label{find:z27}
\COMP{$\mathcal C$} No configuration in $\mathcal C$ pushes the $\zeta_2(7)$ margin
positive. The best margin is $-5.4548$, at $p=2$, $r=1$, $A=4$, $m=3$, $\epsilon=1$, all
shifts $\tfrac12$:
\[
 G=5.5451774444795624753=8\log2,\quad
 \alpha=11.090354888959124951=16\log2=2G,
\]
\[
 \beta=G+11=16.545177444479562475,\quad
 \marg=-5.4548225555204375247 .
\]
The uncorrected ledger would report $\Eled=7$ and $\marg=-1.4548$; the four-unit gap is the
second prime band of Finding~\ref{find:R1R2}, which no shape in the cone removes
(Finding~\ref{find:bandshape}). Re-verified at $50$ digits, and $\Etrue$ re-verified at
larger $n$: the bands are $(7,4)$ at $n=48,60,72$, and the measured $\log(\mathrm{denom})/n$
is $9.72,10.18,10.21$, rising to $11$ from below.
\end{finding}

\subsection{The escape hatch, and why it is shut}\label{ss:hatch}

By Proposition~\ref{prop:pge5} the $p\ge5$ deficit is $-E$ where its hypothesis holds, and
at worst $A\,p\log p/(p-1)^2-E$ in general (Corollary~\ref{cor:noAhelp}); either way the
single inequality in the derivation is $\alpha=\min_j v_p(S_n(j/p^r))$: a twisted sum could in principle be
%% REFEREE [R1]: the unconditional reading of Prop.~\ref{prop:pge5} is repaired here too.
\emph{more} $p$-adically small than its worst summand, through systematic cancellation
between cosets. Nothing in the literature does this and the ledger has no term for it. We
tested it.

Since the partial fractions $r_{i,k}$ do not depend on the integration shift, one exact
computation serves every coset; the $J_u$ come from the exact truncated Volkenborn series.
Two questions were asked.

\smallskip\noindent\textbf{(1) The sums the theory actually needs.} The plain sum and all
$\varphi(p^r)$ Teichm\"uller twists $\sum_j\omega(j)^kS_n(j/p^r)$:
\begin{center}
\tabsmall
\begin{tabular}{lccc}
\toprule
configuration & $\min_jv_p$ at the largest $n$ & best twisted sum & gain\\
\midrule
$p{=}5$, $A{=}2$, bricks on $1/5$, $n{=}27$ & $70$ & $71$ & $+1$\\
$p{=}5$, $A{=}4$, bricks on $1/5$, $n{=}21$ & $109$ & $109$ & $0$\\
$p{=}5$, $A{=}4$, spread over $4$ cosets, $n{=}21$ & $134$ & $136$ & $+2$\\
$p{=}5$, $A{=}4$, $2{+}2$ reflection split, $n{=}21$ & $109$ & $110$ & $+1$\\
$p{=}5$, $A{=}4$, $m{=}1$, $\epsilon{=}1$, $n{=}21$ & $109$ & $110$ & $+1$\\
$p{=}7$, $A{=}3$, $\epsilon{=}1$, on $1/7$, $n{=}18$ & $65$ & $65$ & $0$\\
$p{=}7$, $A{=}3$, spread, $n{=}18$ & $65$ & $66$ & $+1$\\
$p{=}5$, $r{=}2$ ($20$ cosets), $n{=}30$ & $139$ & $140$ & $+1$\\
\bottomrule
\end{tabular}
\end{center}
Gains are $0$ to $3$ in \emph{absolute} terms and do not grow with $n$. The requirement is
$\Omega(n)$; specifically $E\cdot n\ge3n$.

\smallskip\noindent\textbf{(2) The sharp question.} Normalise
$u_n:=(S_n(j))_j/p^{\min_jv_p}$, so that some coordinate is a unit. A \emph{fixed}
$c\in\Zp^{\varphi}$ gains $\Omega(n)$ for all $n$ if and only if all the directions $u_n$ lie
in one hyperplane to depth $\Omega(n)$; that depth is $H:=\max_id_i$ of the Smith normal
form of the matrix $(u_n)$ over $\Zp$ (unit-tested on synthetic data). Measured, once the
number of $n$-values exceeds the number of cosets:
\begin{center}
\tabsmall
\begin{tabular}{ll}
\toprule
configuration & $H$\\
\midrule
$p=5$, $A=2$, on $1/5$ & $14,14,14,14$ at $n=15,18,21,24,27$ (while $\min v_p$ grows $40\to70$)\\
$p=5$, $A=4$, on $1/5$ & $26,26,26$ at $n=15,18,21$\\
$p=5$, $A=4$, spread & $0,0,0$ at $n=15,18,21$\\
$p=5$, $A=4$, $2{+}2$ & $20,20,20$ at $n=15,18,21$\\
$p=5$, $A=4$, $m=1$, $\epsilon=1$ & $28,28,28$ at $n=15,18,21$\\
$p=7$, $A=3$, on $1/7$ & $11,11,11$ at $n=14,16,18$\\
$p=7$, $A=3$, spread & $10,10,10$ at $n=14,16,18$\\
\bottomrule
\end{tabular}
\end{center}
\emph{$H$ is constant in $n$} while $\min_jv_p(S_n)$ grows linearly. Hence no fixed
$\Zp$-combination of the cosets --- twisted, plain or arbitrary --- gains more than a bounded
number of digits.

\begin{finding}[the hatch is shut]\label{find:hatch}
\EXCL{exact $p$-adic arithmetic; tolerance zero (exact valuations); range: $p\in\{5,7\}$,
$r\in\{1,2\}$, $A\in\{2,3,4\}$, $m\in\{0,1\}$, $\epsilon\in\{0,1\}$, bricks aligned / spread
/ reflection-split, $n\le27$ at $p=5$, $n\le18$ at $p=7$, $n\le30$ at $p=5$, $r=2$}
The gain over the per-coset minimum is $\le3$ absolutely, and the hyperplane depth is
$\le28$ and constant in $n$, while the requirement grows like $E\cdot n$.
\end{finding}

\begin{remark}[what this is and is not]\label{rem:hatchcaveat}
This is not a proof of impossibility. It is a fair sweep with a clean negative and a stated
range. One caveat is recorded in place: the $p=5$, $r=2$ run has $20$ cosets and only $9$
values of $n$, so its $H$ is rank-deficient and uninformative; its $\omega$-twist gain is
$+1$, which is the informative half of that row.
\end{remark}

\subsection{Non-vanishing, for the record}\label{ss:nonvanrecord}

No winner exists in $\mathcal C$, so nothing here needs a non-vanishing argument. For the
record, the two objects that would have needed one:
\begin{itemize}[leftmargin=1.6em,itemsep=3pt]
\item the $M=6$, $m=1$ rank-$2$ family (margin $+1.3178$) inherits LSZ's route: the parity
      law gives $\rho_2=\rho_4=\rho_6=0$ identically, and the surviving $\rho_3,\rho_5$ are
      nonzero at every $n$ tested \VER{exact, $n\le60$}; a Casoratian argument in the style
      of \cite[Lemma 15(b)]{LSZ2025} would be the natural proof. It is moot, by
      Finding~\ref{find:rankwall};
\item the $m=-1$ route is dead on smallness, not on vanishing --- though we note the
      degeneracy $\rho_0\equiv0$ identically for $M$ odd with $m=-1$, and for $M$ even with
      $m$ even.
\end{itemize}
Orbit moves preserve non-vanishing (Proposition~\ref{prop:orbitnonvan}); that remains
available and, in the present cone, unused.
